\begin{textsample}{POS Dim 3 – ba – Score 15.00 – ba/ba\_inf\_7.txt}  \label{ex:f3_pos_177}
7.7.
As Transições entre Educação \textbf{Infantil} e Ensino Fundamental : Questões Didáticas e Curriculares É num cenário da Educação \textbf{Infantil} impactado não apenas pelo paradigma da necessidade, mas, a fortiori, pelo paradigma do direito da \textbf{criança} à educação, que a ideia forte hoje em pauta no discurso oficial sobre o “ direito à aprendizagem ” se configura.
Entretanto, da nossa perspectiva, o conteúdo desse discurso precisa de uma reflexão, ou mesmo de uma inflexão mais que fundamental em realidade fundante.
Senão, vejamos.
Para o campo das práticas educacionais, a aprendizagem é uma pauta de interesse inarredável em termos pedagógicos.
Concordamos, entretanto, como vimos afirmando, em termos valorativos nem toda aprendizagem é boa, levando em conta a perspectiva de quem vive a experiência do aprender, suas necessidades, demandas e direitos.
Para a educação, como lugar eminentemente político e ético, a aprendizagem só se legitima se for valorada e referenciada socialmente.
Nesse sentido, o que constitui formação para uma \textbf{criança} vai muito além do que está prescrito, do que deve aprender, nos provoca Pierre Dominicé ( 2012 ).
Sem nos alongar nessa argumentação introdutória, mas importante para explicitarmos nossas intenções de texto, ao se falar em “ direito à aprendizagem ”, fica -se a meio caminho de uma complexidade fundante e de uma necessidade socialmente referenciada : o direito à ( aprendizagem ) formação.
É com esse sentido político-pedagógico do aprender formando -se e do aprender a formar -se na escola de Educação \textbf{Infantil} e das séries iniciais do Ensino Fundamental que caminharemos com nosso texto político-curricular, em que a provocação central é jamais desvincular aprendizagem da formação, envolvendo aí as questões do currículo da educação das \textbf{crianças} \textbf{pequenas}.
Por outra perspectiva mais ampla, Silva, Pasuch e Silva ( 2012 ) explicitam que vivenciamos, nas décadas de 1980 e 1990, a transição da Educação \textbf{Infantil} pautada no paradigma da necessidade para um paradigma do direito da \textbf{criança} sujeito de direitos.
Para a autora, essa transição marca o reconhecimento da importância da Educação \textbf{Infantil} para o processo de formação da \textbf{criança} \textbf{pequena} em ambiências de \textbf{cuidado} e de aprendizagens organizadas para educá -la.
Quanto à didática da Educação \textbf{Infantil}, inspirada no trabalho com os Campos da Experiência, torna -se necessário definir os objetivos centrais da Pedagogia da \textbf{Infância}.
De acordo com Moysés Kuhlmann Jr. ( 1999 ), não se trata de dizer que o conhecimento e a aprendizagem não pertençam ao universo da Educação \textbf{Infantil}, porém a dimensão que assumem na educação das \textbf{crianças} \textbf{pequenas} coloca -os numa relação extremamente vinculada aos processos gerais de constituição da \textbf{criança} : a expressão, o \textbf{afeto}, a sexualidade, a socialização, o \textbf{brincar}, a linguagem, o movimento, a fantasia e o imaginário.
Tal didática não é aquela “ por meio das quais se incentivam aprendizagens particulares ”, mas a que busca aprofundar os tipos de experiência que as \textbf{crianças} vivem e podem viver diariamente [... ] ( a manipulação, o rabisco, os transformismos, o \textbf{contar} histórias, as atividades motoras, o comentário de figuras ) ” e, ao mesmo tempo, tenta “ identificar situações inéditas que possam incentivar nas \textbf{crianças} a exploração e transformação do ambiente ” ( BONDIOLI ; MANTOVANI, 1998, p. 31 ).
Como a realidade da \textbf{criança} é ainda bastante fragmentada, marcada pelo “ aqui e agora ”, a possibilidade de continuidade garante o \textbf{crescimento} e a qualidade das experiências dos \textbf{meninos} e das \textbf{meninas}.
Isso significa dizer que a continuidade implica condições objetivas ( i ) de tempo, para que as \textbf{crianças} possam permanecer em seus percursos de investigação ; ( ii ) de materiais em quantidade suficiente, para que cada \textbf{criança} do grupo não seja constantemente interrompida, e com variedade ampliada para aumentar seu repertório de negociações entre os próprios materiais ; ( iii ) de espaço, pois se faz necessário garantir opções diversas de atuações das \textbf{crianças} em um mesmo local, sem que as obrigue a permanecer todas em uma mesma atividade por longos períodos de tempo ; e ( iv ) de grupo, pois já se sabe que as \textbf{crianças} \textbf{conseguem} atuar melhor quando estão em \textbf{pequenos} grupos.
Novamente, alertamos para o fato de que os campos de experiências não operam em tempos compartimentados : eles atravessam de forma objetiva o modo como o contexto é organizado e, subjetivamente, nas ações e intervenções do \textbf{adulto} que os acompanha.
Em especial, quando nos referimos ao princípio da continuidade, estamos a convocar os professores para o seu papel de responder a essa “ forte exigência de continuidade. [... ] Estabelecendo \textbf{hábitos}, isto é, momentos reconhecíveis pela sua identidade e repetitividade, ou ainda [... ] prestando atenção às possibilidades intrínsecas a cada experiência ” ( BONDIOLI ; MANTOVANI, 1998, p. 32 ).
Na continuidade das experiências é que reside a força e a vitalidade da ação das \textbf{crianças} em compreender, explorar e aprofundar as suas hipóteses afetivas, cognitivas e sociais sobre o mundo.
A partir desse ponto de vista, “ o princípio da continuidade da experiência significa que toda experiência tanto toma algo das experiências passadas quanto modifica de algum modo a qualidade das experiências do presente e que virão ” ( DEWEY, 2010, p. 36 ).
Daí o convite a outro princípio, a significatividade.
O caráter \textbf{lúdico} e contínuo das experiências das \textbf{crianças} abre um espaço para a produção da ludicidade, da continuidade e da significatividade das experiências \textbf{infantis} que podem servir como guias importantes na elaboração e organização dos currículos para dar significados pessoais, seja pelo prazer do já-vivido característico na atividade \textbf{lúdica}, seja por germinar algo que está embrionário na \textbf{criança} na continuidade de suas experiências.
Bondioli e Mantovani ( 1998, p. 32 ) nos indicam estes princípios como sendo “ de uma didática do fazer com as \textbf{crianças} \textbf{pequenas} ” que, aqui neste texto, serviram para ampliar e nos fazer pensar sobre a experiência nos campos de experiência.
Na Pedagogia de Projetos, por exemplo, a abordagem educativa está centrada na \textbf{criança} e na possibilidade de que ela possa expressar -se por meio de todas as suas linguagens, como desenhos, \textbf{pinturas}, palavras, movimentos, dobraduras, modelagens, montagens, dramatizações, colagens, esculturas, músicas, o que a conduz a uma ampla possibilidade de expressão e de criatividade.
A organização das atividades baseia -se numa contínua e responsável flexibilidade e inventividade didática em relação à variabilidade dos ritmos, dos tempos e dos estilos de aprendizagem, além das motivações e dos interesses das \textbf{crianças}.
A organização da programação pedagógica, a partir do trabalho centrado nas experiências, além de valorizar as experiências das \textbf{crianças}, favorece ainda uma organização flexível e autônoma, na qual cada escola de Educação \textbf{Infantil} possa valorizar as suas experiências e culturas locais.
Permite ainda que a escola da \textbf{infância} possa ter flexibilidade e autonomia para construir um projeto próprio, criativo a respeito das orientações das políticas de currículo.

% matched lemmas: adulto, afeto, brincar, conseguir, contar, crescimento, criança, cuidado, hábito, infantil, infância, lúdico, menino, pequeno, pintura
\end{textsample}
